\section{工程实践与控制}

\subsection{Proof Search 调度}
Hubert 介绍了一个调度器：Edges 具有不同 cost（Heuristic, Beam, MCTS），每次在 chunk gating 表中选择 search policy；Chunk gating 记录了 token count、Overlap 以及 proof term size。调度器根据 latency budget 选用“浅 search + high replay”或“深 MCTS + low replay”。
\begin{table}[H]
\centering
\begin{tabular}{p{4cm}p{4cm}p{4cm}}
\toprule
Search 类型 & 特征 & 适用场景 \\
\midrule
Heuristic tactics & 低 latency & 前期 exploration \\
Beam search & 保守选择 top-k & lemma discovery \\
MCTS + RL policy & 状态 value + policy & 最终 polish，适合定理级 \\
\bottomrule
\end{tabular}
\caption{Proof search 调度策略}
\end{table}
\begin{importantbox}{调度器的 KPI}
Latency budget、proof-term length、entropy drift 共同驱动 chunk gating gate；超过阈值就降级到 conservative search。
\end{importantbox}

\subsection{Trace 日志与验证}
每次 failed proof 都会录入 `trace.json`：包含 tactic sequence、Lean 提示、search depth。Hubert 强调 trace 里必须有 `human comment` 字段，使 governance team 能快速 triage。他还指出：Lean 验证失败后，RL 需要 replay alternate candidate 以避免 stuck。
\begin{warningbox}{Trace 日志的警示}
如果没有 replay (trail) 就不可能衡量 strategy drift；没有 trace 的 replay buffer 会让 policy 盲目训练 wrong proofs。
\end{warningbox}

\subsection{人机治理与重现}
治理团队会定期把 high-drift chunk 送到 human-in-the-loop board：reviewer 检查 trace/log、评估 fairness、在 newsletter 中更新 knowledge summary。这种做法让 AlphaProof 将每个失败当作复现实验的 seed。
\begin{knowledgebox}{治理要点}
1）High drift chunk 触发 alert；2）Reviewer 生成 corrected tactic；3）Policy replay + reward correction；4）Newsletter 里记录 QA 亮点。
\end{knowledgebox}

\subsection{本章小结}
工程部分强调 Proof search 调度、trace 记录与治理闭环：通过 chunk gating、trace + review、newsletter 三角形，把 RL agent 从未解释的 black box 变成可复现的研发生命线。
\newpage

